% Validazione dell'ordinamento del divario formativo
% Richiede: \usepackage{booktabs}
\begin{table}[htbp]
  \centering
  \caption{Confronto fra la prossimità fra competenze ricostruita per
           co-occorrenza e la gerarchia ufficiale ESCO v1.2.1. Le coppie
           selezionate dalla ricostruzione sono 2,5 volte più vicine di
           coppie casuali, ma il numero di co-occorrenze non predice la
           vicinanza tassonomica: la ricostruzione conteneva informazione
           binaria e veniva usata come misura graduata.}
  \label{tab:skill-graph-accordo}
  \begin{tabular}{lrr}
    \toprule
    \textbf{Indicatore} & \textbf{Coppie co-occorrenti}
      & \textbf{Coppie casuali} \\
    \midrule
    Coppie esaminate                       & 10.000 & 10.000 \\
    Similarità tassonomica media           & \textbf{0,103} & 0,040 \\
    Quota con antenato comune              & 28,3\% & 13,9\% \\
    Quota concettualmente vicine ($\geq$ 0,5) & \textbf{7,9\%} & 1,8\% \\
    \midrule
    \multicolumn{3}{l}{Correlazione di rango co-occorrenze $\sim$ similarità:
      $\rho = 0{,}033$} \\
    \bottomrule
  \end{tabular}
\end{table}

\begin{table}[htbp]
  \centering
  \caption{Contributo effettivo delle tre componenti del punteggio di
           priorità, misurato su 787 divari formativi reali. Lo scarto fra
           peso nominale e contributo effettivo dipende dalla scala su cui
           ciascuna componente si distribuisce; per la co-occorrenza deriva
           da una normalizzazione che divide una media di co-occorrenze per
           una frequenza massima.}
  \label{tab:skill-graph-contributi}
  \begin{tabular}{lrr}
    \toprule
    \textbf{Componente} & \textbf{Peso nominale}
      & \textbf{Contributo effettivo} \\
    \midrule
    Trasversalità (\texttt{reuse\_level}) & 40\% & \textbf{72,6\%} \\
    Frequenza fra le occupazioni          & 40\% & 20,7\% \\
    Co-occorrenza ricostruita             & 20\% & \textbf{1,6\%} \\
    \midrule
    Prossimità tassonomica (sostituzione) & 20\% & 6,7\% \\
    \bottomrule
  \end{tabular}
\end{table}

\begin{table}[htbp]
  \centering
  \caption{Effetto della sostituzione sull'ordinamento finale del divario, a
           pesi invariati. Il campione è costruito per ciascun curriculum come
           differenza fra le competenze essenziali dell'occupazione di
           riferimento e quelle estratte dal parser: 787 casi su 1.346
           competenze distinte.}
  \label{tab:skill-graph-effetto}
  \begin{tabular}{lr}
    \toprule
    \textbf{Indicatore} & \textbf{Valore} \\
    \midrule
    Casi esaminati                            & 787 \\
    Tau di Kendall medio                      & \textbf{0,791} \\
    Ordinamento invariato                     & 20,5\% \\
    Prima competenza raccomandata cambiata    & \textbf{14,2\%} \\
    \bottomrule
  \end{tabular}
\end{table}
